\section{神经网络为何需要非线性}
\label{sec:14-Deep-Learning/01-Representation-and-Approximation/01}

一个表格数据的回归任务卡住了。线性模型在验证集上的决定系数停在 \(0.62\)，怎么调正则强度都推不动。有人提了个听上去很合理的方案：既然深度学习厉害，那就把这一层线性映射拆成三层堆起来，参数量翻两番，总该有点长进。改完，训练跑通，验证集上的数字一动不动——不是波动了一点点，是几乎逐位相同。学习率、初始化、批大小逐一排查，全都正常。

问题不在训练过程，在于改完之后的模型类和改之前是同一个。这一节要说清楚这次塌缩为什么必然发生，以及插进一个激活函数究竟改变了什么。

\subsection{线性层叠起来仍然是一个线性层}

把两层写出来：

\[
h_1=W_1x+b_1,\qquad
h_2=W_2h_1+b_2.
\]

代入展开，

\[
h_2=(W_2W_1)x+(W_2b_1+b_2).
\]

右端是一个矩阵乘 \(x\) 再加一个向量，形状和单层完全一致。继续往下叠 \(L\) 层，结果仍是 \(W_{L}\cdots W_1\) 这一个矩阵配一个偏置。深度在这里是白加的：参数变多了，可表示的函数集合一点没变。

这不是神经网络才有的怪事，而是线性映射的基本性质——线性映射的复合仍是线性映射，矩阵乘法正是这条性质的坐标表达，\hyperref[sec:03-Linear-Algebra/02-Vector-Spaces/04]{``线性变换与换基：对象不变，坐标可以改变''}把这件事讲得更细。所以那次改造真正做的事情，是把同一个映射换了一种参数化写法。原来用 \(d_{\text{out}}\times d_{\text{in}}\) 个自由参数描述的东西，现在用三个矩阵的乘积描述，秩还可能因为中间层变窄而被压低。函数类没扩大，反而可能缩小。

顺带一提，若中间层维度小于输入输出维度，乘积矩阵的秩被中间层卡住，这时候的``深度''等价于一次低秩约束。它有用，但用处属于正则化，不属于表达能力。

\subsection{激活函数的作用是把这次塌缩挡住}

在两次仿射变换之间插入一个逐元素的非线性函数 \(\phi\)：

\[
h_2=W_2\,\phi(W_1x+b_1)+b_2 .
\]

只要 \(\phi\) 不是仿射函数，乘积就无法再合并。以 ReLU 为例，

\[
\relu(z)=\max(0,z),
\]

单个 ReLU 单元给出一条折线：阈值左侧压平为零，右侧按斜率上升，折点位置由偏置决定。折点是线性映射里造不出来的东西——线性函数处处一个斜率。

三个折线叠起来就能造出局部结构。取

\[
f(x)=\relu(x)-2\relu(x-1)+\relu(x-2),
\]

在 \(x<0\) 上为零，在 \((0,1)\) 上以斜率 \(1\) 上升，在 \((1,2)\) 上以斜率 \(-1\) 下降，越过 \(2\) 之后重新归零。三个单元拼出了一个三角形的局部凸起，而它在整条实轴上不是任何一条直线。

\begin{center}
\begin{tikzpicture}[>=stealth,scale=1.0]

\draw[->] (-0.5,0) -- (3.5,0) node[right,font=\small] {\(x\)};
\draw[->] (0,-0.5) -- (0,2.9) node[above,font=\small] {输出};
\draw[very thick] (-0.45,0.05) -- (3.3,2.75);
\node[font=\small] at (2.05,0.75) {\(W_2W_1x+b\)};
\node[font=\small] at (1.5,-1.05) {叠多少层都还是它};

\draw[->,very thick] (4.0,1.3) -- (5.2,1.3);
\node[font=\small] at (4.6,1.65) {插入 \(\relu\)};

\begin{scope}[shift={(5.9,0)}]
  \draw[->] (-0.5,0) -- (3.5,0) node[right,font=\small] {\(x\)};
  \draw[->] (0,-0.5) -- (0,2.9) node[above,font=\small] {输出};
  \draw[dashed,gray] (1,0) -- (1,1.8);
  \draw[dashed,gray] (2,0) -- (2,0.15);
  \draw[very thick] (-0.45,0) -- (0,0) -- (1,1.8) -- (2,0) -- (3.3,0);
  \fill (0,0) circle (2pt);
  \fill (1,1.8) circle (2pt);
  \fill (2,0) circle (2pt);
  \node[font=\small] at (1.5,-1.05) {三个折点，三段不同斜率};
\end{scope}

\end{tikzpicture}
\end{center}

\emph{左边无论叠多少层都只是一条直线，右边三个 hinge 相加已经折出一个局部三角。}

图里右侧那条折线是理解整个深度模型的最小单位。网络做的事情，是把大量这样的折点摆到高维空间里，再让后面的层在折出来的分区上继续折。

\subsection{一层里同时发生了切割和折叠}

前馈网络的标准写法是

\[
h_0=x,\qquad
a_\ell=W_\ell h_{\ell-1}+b_\ell,\qquad
h_\ell=\phi_\ell(a_\ell).
\]

仿射部分负责重新组合坐标、平移分界面，把输入投到一组新的方向上；激活函数负责在这些方向上截断或弯折。对 ReLU 网络，每个单元把空间切成``激活''与``不激活''两个半空间，所有单元的激活模式共同把输入空间划分成许多多面体区域。在每个区域内部，整个网络退化成一个仿射映射；区域与区域之间，斜率不同。

深度带来的差别在于折叠是复合的。第一层的折点画在输入空间上，第二层的折点画在第一层输出的坐标里——同一个第二层单元，可能在输入空间中对应好几段互不相邻的边界。前一个三角形若被再次送进类似的构造，折点数量会随层数成倍增长，而单层网络要摆出同样多的折点，只能靠把单元数堆到同一个量级。这条线索下一节会正式接手。

\subsection{深度模型把选表示这件事变成了可优化的对象}

值得停下来对比一下机器学习那一部分的做法。核方法里，特征映射 \(\varphi\) 由人挑定：多项式核、RBF 核、字符串核，选哪一个是建模者的判断，选定之后模型只在这个固定的特征空间里求一个线性解——\hyperref[sec:13-Machine-Learning/09-Kernel-Methods-and-Gaussian-Processes/01]{``把数据升到高维为什么有用''}把这条思路交代得很完整。树模型的分裂规则、线性模型的交互项、时间序列的滞后阶数，同样都由人事先写死。表示的好坏在训练开始之前就已经决定了大半。

带非线性的多层网络改掉的正是这一点。\(W_1,b_1\) 既是参数也是特征映射的定义，梯度既更新分类边界也更新用来分类的那组特征。选表示从建模者的一次判断，变成了目标函数上的一个可微变量。

代价来得同样直接。核方法给出的解是凸问题的全局最优，Gaussian process 甚至给得出后验方差；一旦特征映射自身进入优化变量，目标函数不再凸，解的存在唯一性、收敛速度、泛化界通通降级为条件苛刻的定性结论或纯粹的经验规律。这是本部分反复出现的一条主线：把决策权从人手里交给梯度，换来的表达力是实打实的，丢掉的保证也是实打实的。下一节讨论的通用逼近定理，是这条主线上第一个也是最著名的例子——它承诺得很响亮，能兑现的部分却窄得多。

\subsection{激活函数同时决定了梯度能不能流回去}

选哪个 \(\phi\) 不只影响能表示什么，还影响训练时信号怎样传回前面的层。sigmoid

\[
\sigma(z)=\frac{1}{1+e^{-z}}
\]

把实数压进 \((0,1)\)，导数是

\[
\sigma'(z)=\sigma(z)\bigl(1-\sigma(z)\bigr)\le\frac14 ,
\]

上界来自 \(s(1-s)\) 在 \(s=1/2\) 处取最大值。\(\abs{z}\) 一大，输出贴近端点，导数趋近于零。反向传播要把这些因子连乘起来，十几层之后传到前端的梯度可以小到浮点数分辨不出来。tanh 以零为中心，缓解了输出偏移，饱和问题照旧。

ReLU 在正半轴导数恒为 \(1\)，连乘不再自动衰减，这是它取代 sigmoid 成为隐藏层默认选择的主要原因。它的麻烦在负半轴：导数为零，某个单元若因为一次大步长被推到对全部样本都不激活的位置，梯度就再也回不来，这个单元事实上死掉了。LeakyReLU 给负半轴留一点斜率，GELU 用光滑过渡替代硬折角，都是在同一处做修补。哪一个更好取决于初始化、归一化和任务本身，没有一个脱离上下文的最优选择。

\subsection{输出层的非线性和隐藏层的非线性不是一回事}

隐藏层的 \(\phi\) 用来制造弯折，输出层的映射用来把网络的实数输出对齐到任务要求的取值范围和概率语义，两者共用``激活函数''这个名字，职责并不相同。无界的连续回归通常直接用线性输出；二元事件用一个 logit 配数值稳定的二元交叉熵；互斥多分类用一组 logits 配 \(\softmax\)；多标签任务要用多个独立 logit，因为标签之间不该在一个 \(\softmax\) 里互相挤压概率质量。这几种选择对应的是不同的概率模型，\hyperref[sec:13-Machine-Learning/02-Linear-Models/03]{``从二分类到 Softmax 分类''}已经把二元与多元的分界讲清楚，下一单元会把``换输出层就是换噪声假设''这件事完整展开。

把隐藏层的激活函数随手搬到输出层，或者反过来在隐藏层用 \(\softmax\)，通常不会报错，只会让训练目标悄悄变成另一个问题。

\subsection{同一个函数可以由许多相距很远的参数表示}

还有一件与非线性直接相关的事：参数空间存在大量对称性。把同一隐藏层的两个单元交换编号，同时交换下一层对应的两列权重，网络实现的函数逐点不变。对 ReLU 单元，把入边权重与偏置乘以 \(c>0\)、出边权重除以 \(c\)，由正齐次性 \(\relu(cz)=c\relu(z)\)，函数同样不变。奇函数激活还额外允许整组符号翻转。

于是两组数值上相距很远的参数可能是同一个函数，两次训练得到的权重矩阵逐元素对比几乎没有意义。要比较模型，得比较输出、比较中间表示张成的子空间，或者比较在受控输入上的行为。这条约束在解释性和模型合并这类工作里会反复出现。

非线性买到的是弯折能力，也一并买来了非凸的参数空间和成堆的等价解。下一节先把``弯折能力到底有多强''这个问题问到底：通用逼近定理承诺一层就够，而这句承诺的适用范围比它的名气小得多。
